% R5 — calibration and robustness axes (draft 2026-09-07)
\subsection{Calibration favors STATE, and random forest probabilities are better calibrated than logistic-regression scores}
\label{ssec:r5_calibration}

Discrimination is not the only property that decides whether a risk estimate
can be acted on: the emitted probability must also mean what it says. We
computed patient-level calibration on the within-COVID out-of-fold predictions
(Figure~\ref{fig:calibration}). Among logistic-regression classifiers,
\model{STATE} was the best calibrated (Brier 0.118, expected calibration error
0.168), ahead of \model{Geneformer} (0.119, 0.190) and \model{scPRINT-2}
(0.121, 0.181), with \model{UCE} (0.162, 0.243) and \model{scGPT} (0.165,
0.257) the least calibrated. Random-forest probabilities were better
calibrated than logistic-regression scores for every model (ECE 0.054--0.080
versus 0.168--0.257), reflecting the conservatism of tree-averaged
probabilities, so where calibrated risk estimates matter for deployment we
recommend random forest on the leading embedding (\model{STATE}-RF: Brier
0.095, ECE 0.079) over any logistic-regression variant.

Two further secondary axes complete the multi-axis profile. First,
\model{UCE} was the most resistant to embedding-dimension dropout: its AUC did
not degrade under 70\% dimension zeroing, indicating a highly distributed
representation (Supplementary Figure~\ref{fig:supp_dropout}; computed on the
legacy analysis pool, an earlier cohort assembly predating the canonical
240-patient cohort, Section~\ref{ssec:limitations}). Second,
leave-one-cohort-out evaluation on the analysis cohort measured
generalization to cohorts that contain fatal cases, an outcome-selected
estimand: folds are computable only when the held-out cohort contains both
classes, so survivor-only cohorts contribute to training throughout and are
never scored out-of-fold (Supplementary Figure~\ref{fig:supp_lobo}).
Because the five computable folds range from 13 to 90 patients and from 1 to
9 fatal events, we report the AUC over concatenated held-out predictions
(206 patients, 25 fatal events) rather than a mean of per-fold AUCs: the top
three models sit within 0.005 AUC of each other (\model{STATE} 0.753,
\model{scPRINT-2} 0.751, \model{Geneformer} 0.748), ahead of \model{UCE}
(0.681) and \model{scGPT} (0.638). Per-fold AUCs and the per-cohort counts
are given in Supplementary Table~\ref{tab:lobo_folds}. By fold-to-fold spread,
\model{scGPT} was the most cohort-consistent ($\sigma = 0.123$) and
\model{UCE} the least ($\sigma = 0.152$). No single model leads on every
axis; the choice among them is deployment-specific
(Section~\ref{ssec:discussion_guidance}).
